\subsection{乘积测度的性质}

若 \(\lambda\) （相应地 \(\mu\) ）是 X（相应地 Y）上的一个测度，而 \(\nu\) 是 \(\lambda\) 与 \(\mu\) 的乘积测度，则 \(\nu\) 在 \(\mathcal{K}(X;C)\otimes_{\mathbf{C}}\mathcal{K}(Y;C)\) 上的限制不是别的，正是张量积 \(\lambda\otimes\mu\) ，即两个线性形式 \(\lambda\) 与 \(\mu\) 的张量积（A, II, §3, No.~2），因为 No.~1 的关系式 (1) 可写成

\[
\langle g \otimes h, \nu \rangle = \langle g, \lambda \rangle \langle h, \mu \rangle = \langle g \otimes h, \lambda \otimes \mu \rangle
\]

对所有 \(g \in \mathcal{K}(X; \mathbf{C})\) 及 \(h \in \mathcal{K}(Y; \mathbf{C})\) 成立。出于记号的滥用，我们也用 \(\lambda \otimes \mu\) 表示乘积测度 \(\nu\) （而不仅是它在 \(\mathcal{K}(X; \mathbf{C}) \otimes_{\mathbf{C}} \mathcal{K}(Y; \mathbf{C})\) （ \(\mathcal{K}(X \times Y; \mathbf{C})\) 的稠密子空间）上的限制）。

映射 \((\lambda, \mu) \mapsto \lambda \otimes \mu\) （ \(\mathcal{M}(\mathrm{X}; \mathbf{C}) \times \mathcal{M}(\mathrm{Y}; \mathbf{C})\) 到 \(\mathcal{M}(\mathrm{X} \times \mathrm{Y}; \mathbf{C})\) 中的）显然是双线性的，这依据 No.~1 的公式 (3)。

命题 1. --- 设 \(\lambda\) 是 X 上的一个测度， \(\mu\) 是 Y 上的一个测度；若 \(g \in \mathcal{C}(\mathrm{X};\mathbf{C})\) ， \(h \in \mathcal{C}(\mathrm{Y};\mathbf{C})\) ，则

\[
(g \cdot \lambda) \otimes (h \cdot \mu) = (g \otimes h) \cdot (\lambda \otimes \mu).\tag{6}
\]

对每个函数 \(f \in \mathcal{K}(\mathrm{X} \times \mathrm{Y}; \mathbf{C})\) ，依据 No.~1 的公式 (3)，有

\[
\begin{array}{c} \langle f, (g \otimes h) \cdot (\lambda \otimes \mu) \rangle = \int d \lambda (x) \int f (x, y) g (x) h (y) d \mu (y) \\ = \int g (x) d \lambda (x) \int f (x, y) h (y) d \mu (y), \end{array}
\]

这就证明了公式 (6)。

命题 2. --- 乘积 \(\lambda \otimes \mu\) 的支集等于 \(\lambda\) 的支集与 \(\mu\) 的支集之积。

我们首先注意，关系式 \(\lambda \otimes \mu = 0\) 蕴涵测度 \(\lambda, \mu\) 之一为零（A, II, §7, No.~7, 命题 16, (ii)）。另一方面，若 U（相应地 V）是 X（相应地 Y）中的开集，则 \(\lambda\otimes\mu\) 在乘积 \(U\times V\) 上的限制是 \(\lambda\) 在 U 上的限制与 \(\mu\) 在 V 上的限制之积，这由 No.~1 的 Th. 1 及测度在开集上的限制的定义（§2, No.~1）得出。于是由此可知，为使 \(\lambda\otimes\mu\) 在 \(U\times V\) 上的限制为零，必须且只须 \(\lambda\) 在 U 上的限制或 \(\mu\) 在 V 上的限制为零，再考虑到测度支集的定义（§2, No.~2），这就证明了该命题。

命题 3. --- 设 \(\lambda \in \mathcal{M}(\mathrm{X};\mathbf{C})\) ， \(\mu \in \mathcal{M}(\mathrm{Y};\mathbf{C})\) 。则

\[
| \lambda \otimes \mu | = | \lambda | \otimes | \mu |.\tag{7}
\]

设 \(f \in \mathcal{K}_{+}(\mathrm{X} \times \mathrm{Y})\) ， \(g \in \mathcal{K}(\mathrm{X} \times \mathrm{Y}; \mathbf{C})\) 满足 \(|g| \leqslant f\) ；我们有（§1, No.~6, 公式 (13)）

\[
\begin{array}{r l} | \langle g, \lambda \otimes \mu \rangle | = & \left| \int d \lambda (x) \int g (x, y)   d \mu (y) \right| \\ & \leqslant \int d | \lambda | (x) \int | g (x, y) |   d | \mu | (y) \\ & = \langle | g |, | \lambda | \otimes | \mu | \rangle \leqslant \langle f, | \lambda | \otimes | \mu | \rangle . \end{array}
\]

由此得 \(\langle f, |\lambda \otimes \mu| \rangle \leqslant \langle f, |\lambda| \otimes |\mu| \rangle\) ，于是

\[
| \lambda \otimes \mu | \leqslant | \lambda | \otimes | \mu |.\tag{8}
\]

另一方面，设 \(u \in \mathcal{K}_{+}(\mathrm{X})\) ， \(v \in \mathcal{K}_{+}(\mathrm{Y})\) 。对每个 \(\varepsilon > 0\) ，存在 \(u_{1} \in \mathcal{K}(\mathrm{X};\mathbf{C})\) ， \(v_{1} \in \mathcal{K}(\mathrm{Y};\mathbf{C})\) ，使 \(|u_{1}| \leqslant u\) ， \(|v_{1}| \leqslant v\) ，且

\[
| \langle u _ {1}, \lambda \rangle | \geqslant \langle u, | \lambda | \rangle - \varepsilon , \quad | \langle v _ {1}, \mu \rangle | \geqslant \langle v, | \mu | \rangle - \varepsilon
\]

（§1, No.~6）。由此得 \(|u_1 \otimes v_1| \leqslant u \otimes v\) ，且

\[
\begin{array}{c} \langle u \otimes v, | \lambda \otimes \mu | \rangle \geqslant | \langle u _ {1} \otimes v _ {1}, \lambda \otimes \mu \rangle | = | \langle u _ {1}, \lambda \rangle \langle v _ {1}, \mu \rangle | \\ \geqslant (\langle u, | \lambda | \rangle - \varepsilon) (\langle v, | \mu | \rangle - \varepsilon). \end{array}
\]

由于 \(\varepsilon\) 是任意的，我们推出

\[
\langle u \otimes v, | \lambda \otimes \mu | \rangle \geqslant \langle u, | \lambda | \rangle \langle v, | \mu | \rangle = \langle u \otimes v, | \lambda | \otimes | \mu | \rangle .
\]

把 (8) 考虑在内，我们便得

\[
\langle u \otimes v, | \lambda \otimes \mu | \rangle = \langle u \otimes v, | \lambda | \otimes | \mu | \rangle .
\]

由于 \(\mathcal{K}(\mathrm{X};\mathbf{C})\) （相应地 \(\mathcal{K}(\mathrm{Y};\mathbf{C})\) ）中的每个函数都是 \(\mathcal{K}_{+}(\mathrm{X})\) （相应地 \(\mathcal{K}_{+}(\mathrm{Y})\) ）中函数的线性组合，前面的公式对 \(u\in\mathcal{K}(\mathrm{X};\mathbf{C})\) 与 \(v\in\mathcal{K}(\mathrm{Y};\mathbf{C})\) 仍成立；于是该命题由 \(\mathcal{K}(\mathrm{X};\mathbf{C})\otimes_{\mathbf{C}}\mathcal{K}(\mathrm{Y};\mathbf{C})\) 在 \(\mathcal{K}(\mathrm{X}\times\mathrm{Y};\mathbf{C})\) 中稠密这一事实得出。

推论. --- 设 \(\lambda \in \mathcal{M}(\mathrm{X};\mathbf{R})\) ， \(\mu \in \mathcal{M}(\mathrm{Y};\mathbf{R})\) 。则

\[
\left\{ \begin{array}{l} (\lambda \otimes \mu) ^ {+} = \lambda^ {+} \otimes \mu^ {+} + \lambda^ {-} \otimes \mu^ {-}, \\ (\lambda \otimes \mu) ^ {-} = \lambda^ {+} \otimes \mu^ {-} + \lambda^ {-} \otimes \mu^ {+}. \end{array} \right.\tag{9}
\]

因为依据命题 3，

\[
\begin{array}{r l} & {(\lambda \otimes \mu) ^ {+} = \frac {1}{2} \big (\lambda \otimes \mu + | \lambda | \otimes | \mu | \big)} \\ & {\qquad = \frac {1}{2} \big ((\lambda^ {+} - \lambda^ {-}) \otimes (\mu^ {+} - \mu^ {-}) + (\lambda^ {+} + \lambda^ {-}) \otimes (\mu^ {+} + \mu^ {-}) \big)} \\ & {\qquad = \lambda^ {+} \otimes \mu^ {+} + \lambda^ {-} \otimes \mu^ {-}.} \end{array}
\]

对 \((\lambda \otimes \mu)^{-}\) 的论证是类似的。

命题 4. --- 设 \(\lambda \in \mathcal{M}(\mathrm{X};\mathbf{C})\) ， \(\mu \in \mathcal{M}(\mathrm{Y};\mathbf{C})\) 。则

\[
\| \lambda \otimes \mu \| = \| \lambda \| \cdot \| \mu \|,\tag{10}
\]

并约定每当一个因子为 0 而另一因子为 +∞ 时，第二成员换成 0。特别地，若 λ 与 μ 有界，则 λ⊗μ 有界。

由上面的命题 3 及 §1, No.~8, 命题 10 的推论 1，可以只限于 \(\lambda\) 与 \(\mu\) 是正测度的情形。若 \(\lambda = 0\) 或 \(\mu = 0\) ，结果是平凡的；故设 \(\lambda \neq 0\) 且 \(\mu \neq 0\) 。先证明 \(\| \lambda \otimes \mu \| \leqslant \| \lambda \| \cdot \| \mu \|\) 。可设 \(\lambda\) 与 \(\mu\) 有界。对每个 \(f \in \mathcal{K}_+(\mathrm{X} \times \mathrm{Y})\) ，

\[
\langle f, \lambda \otimes \mu \rangle = \int d \lambda (x) \int f (x, y) d \mu (y)
\]

且

\[
\int f (x, y) d \mu (y) \leqslant \| f \| \cdot \| \mu \|
\]

对每个 \(x\in \mathbf{X}\) 成立，因此

\[
\langle f, \lambda \otimes \mu \rangle \leqslant \| f \| \cdot \| \lambda \| \cdot \| \mu \|,
\]

这就证明了我们的断言。另一方面，设

\[
\alpha <   \| \lambda \|, \beta <   \| \mu \|
\]

为两个实数 \(\geqslant 0\) 。存在 \(g\in \mathcal{K}_{+}(\mathrm{X})\) ， \(h\in \mathcal{K}_{+}(\mathrm{Y})\) ，使

\[
\| g \| \leqslant 1, \| h \| \leqslant 1, \lambda (g) \geqslant \alpha , \mu (h) \geqslant \beta .
\]

于是 \(g \otimes h \in \mathcal{K}_{+}(X \times Y)\) ， \(\| g \otimes h \| \leqslant 1\) ，且 \(\langle g \otimes h, \lambda \otimes \mu \rangle \geqslant \alpha \beta\) ；因此 \(\| \lambda \otimes \mu \| \geqslant \alpha \beta\) ，最终得 \(\| \lambda \otimes \mu \| \geqslant \| \lambda \| \cdot \| \mu \|\) ，这就完成了证明。

